Putting
@@M0@@, we are reduced to calculating the residue of:
@@M1@@
that is, the coefficient of @@M2@@ in @@M3@@. For @@M4@@, it is zero, and for @@M5@@, it is @@M6@@.

This completes the proof of the Riemann--Roch theorem.

\sgaStableTarget{sga6.e08.bibliography}\subsection*{BIBLIOGRAPHY}
\begin{enumerate}[label={[\arabic*]}]
\item A. Borel and J.P. Serre, Le théorème de Riemann--Roch, Bull. Soc. Math. de France, t. 86, 1958.
\item A. Grothendieck, Techniques de construction et théorèmes d'existence en Géométrie Algébrique, Sém. Bourbaki, 1960-61.
\end{enumerate}
